% ============================================================
% 注意事項（手引書のはじめに入れる）
% 原本: 電気電子工学実験I_配布資料260927.pdf の「注意事項」
% ============================================================

\introtitle{注意事項}

\begin{enumerate}[label=\textbf{\arabic*.}, leftmargin=*, labelsep=0.8em, itemsep=0.6\baselineskip]
  \item あらかじめ手引書を熟読し，内容をよく理解して実験に臨んでください。
  \item 毎回遅刻することなく履修してください。遅刻は減点の対象となります。また，\wave{長時間（概ね10分以上）の遅刻は欠席とみなします。}
  \item やむをえない理由で遅刻，欠席する場合は，必ず事前に担当教員に連絡し，許可を得てください。\wave{無断欠席は以降の履修を認めません。}何らかの事情で連絡が取れない場合でも，早急に担当教員に連絡を取ってください。
  \item \wave{電気を利用する実験では，水に濡れている状態は非常に危険です。}事故を回避するため，実験中に誤っても水に濡れることがないように，以下の点に注意して実験を行ってください。
    \begin{itemize}
      \item 濡れた傘は実験室に持ち込まないでください。
      \item 実験中に水分補給が必要な際には，実験室の外に出て飲んでください。
      \item 実験室ではペットボトルや水筒等をかばんの中にしまっておいてください。
    \end{itemize}
    また，実験ごとの留意点は担当教員から注意があるので，実験に当たっては担当教員の指示に従ってください。
  \item 報告書（レポート）は担当教員から指示される期限までに提出してください。遅れた場合は遅れの期間に応じて減点されます。ただし，やむをえない理由で前もって申し出た場合，提出期限の延長を許可されることがあります。
  \item 報告書の修正等の指示がなされた場合，指示された期限までに再提出し，受理されるようにしてください。期限までに再提出されない場合は，報告書未提出とみなすことがあります。また，\wave{報告書が1通でも未提出の場合は単位を取得できません。}
  \item 再提出の機会は原則1回のみで，その状態で報告書は採点・評価されます。
\end{enumerate}
